As simulações foram realizadas no ambiente Simulink utilizando a planta
identificada de acordo com a \Eqref{planta-sensor-1-zn} e os controladores das
\Eqref{controlador-pid-paralelo,controlador-ff}, discretizados pelo método de
Tustin com período de amostragem de $T_s=\SI{15}{\second}$. Para aproximar as
condições experimentais, foi adicionado ruído branco à saída da planta,
representando perturbações e incertezas de medição.

Os resultados obtidos demonstram que a estratégia composta por \gls{pid} e
\gls{ff} foi capaz de acompanhar a referência com erro reduzido ao longo do
transitório e em regime permanente.